% SHARD-ID: CA-01-PROGRAMME-MAP
% SHARD-TITLE: Programme Map
% SHARD-SUMMARY: Lays out the construction pipeline as questions: category data, microscopic models, lattice symmetry generators, provable continuum limit, rigorous CFT.
% SHARD-SUMMARY: Records the first convention and source decisions needed before technical derivations begin.
% SHARD-KEYWORDS: programme, fusion category, modular tensor category, anyon chain, string-net, Koo-Saleur, continuum limit, conformal net

\section{Programme Map}

This shard is a map of questions, not a list of results. The north star
(\S\ref{sec:frontmatter} / CA-00) factors into a pipeline of stages; each stage
is an open problem with its own conventions, sources, and checkable invariants.

\subsection{The Pipeline}

\[
  \underbrace{\Cc \;(+\,\text{OPE/conformal data})}_{\text{input}}
  \;\longrightarrow\;
  \underbrace{\text{microscopic models}}_{\text{lattice Hilbert space + Hamiltonian}}
  \;\longrightarrow\;
  \underbrace{\text{lattice symmetry generators}}_{\text{e.g.\ Koo--Saleur } L_n}
  \;\longrightarrow\;
  \underbrace{\text{provable continuum limit}}_{\text{OAR / wavelets}}
  \;\longrightarrow\;
  \underbrace{\text{rigorous (C)FT}}_{\text{conformal net / VOA / Wightman}}
\]

\subsection{Stage 1 --- Category Input}

\begin{itemize}
  \item Which categories are in initial scope: multiplicity-free unitary fusion
  categories (Fibonacci, Ising/$\mathrm{TL}$), then modular tensor categories?
  \item What auxiliary data beyond $\Cc$ is genuinely required (OPE coefficients,
  central charge, conformal weights), and for which target CFTs?
  \item Which gauge and indexing conventions are fixed before any F/R-symbol is
  written (see \texttt{CONVENTIONS.md})?
\end{itemize}

\subsection{Stage 2 --- Microscopic Models}

\begin{itemize}
  \item Which family realises $\Cc$ on the lattice: anyon chains, string-nets,
  or another construction --- and how is the family indexed by scale?
  \item What is the lattice Hilbert space (fusion-tree basis) and the
  Hamiltonian, and which invariants pin them (fusion-rule consistency, pentagon,
  projector idempotence)?
\end{itemize}

\subsection{Stage 3 --- Lattice Symmetry Generators}

\begin{itemize}
  \item Which lattice operators approximate the continuum symmetry generators
  (Koo--Saleur lattice Virasoro $L_n$, $\bar L_n$, and kin)?
  \item What convergence statement is the target, and what is checkable at finite
  size (commutator algebra, central charge extraction)?
\end{itemize}

\subsection{Stage 4 --- Provable Continuum Limit}

\begin{itemize}
  \item Which rigorous renormalisation framework carries the limit (operator-
  algebraic renormalisation, wavelet / multi-scale entanglement)?
  \item What must be proved: existence of the limit, the limiting Hilbert space,
  and convergence of the symmetry generators to a projective unitary
  representation?
\end{itemize}

\subsection{Stage 5 --- Rigorous (C)FT Target}

\begin{itemize}
  \item What is the target object: a conformal net, a vertex operator algebra, or
  a Wightman/Osterwalder--Schrader theory?
  \item What is the fidelity statement: that the target's superselection /
  representation category recovers $\Cc$?
\end{itemize}

\subsection{First Decisions}

The first technical session should not start with formulas. It should register
the sources for one worked input category (Fibonacci or Ising/$\mathrm{TL}$),
fix the minimum conventions (F/R-symbol gauge, vacuum index, fusion-tree
ordering), and identify the smallest invariant that can be checked end-to-end on
one stage of the pipeline.
